% !TEX root = ../../../main.tex
% 所属：第二部分：多变量分析学 / 第3章：重积分
% 节：3.4 反常重积分
% 正文按原文件顺序保存；由 ../chapter3.tex 调用。
\section{反常重积分}

两类：一是无穷区域，二是有界区域上的无界函数。

\subsection*{无穷区域上的反常二重积分}

设 $D\subset\R^2$ 为无穷区域，其边界由有限条光滑曲线组成（或没有边界）。
被积函数 $f(x,y)$ 在 $D$ 中有界，在 $D$ 中任意一个有界区域内都有通常意义下的二重积分。
从 $D$ 中用一条零面积闭曲线 $\Gamma$ 截取一个有界区域，记为 $D_\Gamma\subset D$。
这里取曲线截取的内部部分，截断曲线向外推移时，区域随之扩大。
记 $\Gamma$ 到原点的距离为
\[
d(\Gamma)=\inf\{\sqrt{x^2+y^2}:(x,y)\in\Gamma\}.
\]
当 $d(\Gamma)\to+\infty$ 时，若
\[
\displaystyle\iint_{D_\Gamma}f(x,y)\,\dd x\,\dd y
\]
有一个与 $\Gamma$ 无关的有限极限，则称 $f(x,y)$ 在 $D$ 上的反常二重积分收敛，
或称 $f(x,y)$ 在 $D$ 上可积，并记
\[
I=\displaystyle\iint_Df(x,y)\,\dd x\,\dd y.
\]
否则称 $f(x,y)$ 在 $D$ 上的反常二重积分发散。
其精确的 $\varepsilon$--$M$ 表述如下。

\begin{definition}[无穷区域上的反常二重积分]
若存在有限数 $I$，使任给 $\varepsilon>0$，都存在 $M>0$，
对所有满足 $d(\Gamma)>M$ 的上述截断曲线，都有
\[
\left|\displaystyle\iint_{D_\Gamma}f(x,y)\,\dd x\,\dd y-I\right|<\varepsilon,
\]
则称 $f$ 在 $D$ 上的反常二重积分收敛，或称 $f$ 在 $D$ 上可积，记为
\[
\displaystyle\iint_Df(x,y)\,\dd x\,\dd y
=\lim_{d(\Gamma)\to+\infty}\displaystyle\iint_{D_\Gamma}f(x,y)\,\dd x\,\dd y=I.
\]
若不存在这样的有限极限，则称反常二重积分发散。
\end{definition}

这里的极限必须与截断曲线的选择无关；仅沿某一种截域算出极限，
对一般的变号函数还不能据此断定反常重积分收敛。

\begin{lemma}\label{p2-improper-exhaustion}
设 $f(x,y)$ 为无界区域 $D$ 上的非负函数。
取一列曲线 $\Gamma_n$，它们割出的 $D$ 的有界子区域 $D_n$ 满足
\[
D_1\subset D_2\subset\cdots\subset D_n\subset\cdots,
\qquad\lim_{n\to\infty}d(\Gamma_n)=+\infty.
\]
则反常积分 $\displaystyle\iint_Df(x,y)\,\dd x\,\dd y$ 在 $D$ 上收敛的充分必要条件是数列
\[
I_n=\displaystyle\iint_{D_n}f(x,y)\,\dd x\,\dd y
\]
收敛。在收敛时成立
\[
\displaystyle\iint_Df(x,y)\,\dd x\,\dd y
=\lim_{n\to\infty}\displaystyle\iint_{D_n}f(x,y)\,\dd x\,\dd y.
\]
\end{lemma}
\begin{proof}
由 $f\ge0$ 及 $D_n\subset D_{n+1}$，$I_n$ 单调不减。

\textbf{必要性。}若 $f$ 在 $D$ 上的反常积分收敛，
那么对任意 $d(\Gamma)\to+\infty$ 的截断曲线，相应积分均收敛于同一数。
特别取 $\Gamma=\Gamma_n$，便得 $I_n$ 收敛。

\textbf{充分性。}设 $I_n\to I$。任取一条截断曲线 $\Gamma$，
它所围出的 $D_\Gamma$ 是有界区域，因此可取
\[
\rho=\sup\{\sqrt{x^2+y^2}:(x,y)\in D_\Gamma\}<+\infty.
\]
由 $d(\Gamma_n)\to+\infty$，存在 $n_0$，使 $n\ge n_0$ 时 $d(\Gamma_n)>\rho$，
因而 $D_n\supset D_\Gamma$。由 $f\ge0$，有
\[
I_n=\displaystyle\iint_{D_n}f(x,y)\,\dd x\,\dd y
\ge\displaystyle\iint_{D_\Gamma}f(x,y)\,\dd x\,\dd y.
\]
令 $n\to\infty$，得到
\[
I\ge\displaystyle\iint_{D_\Gamma}f(x,y)\,\dd x\,\dd y.
\]

另一方面，由 $I_n\to I$，对任意 $\varepsilon>0$，存在 $N$，使 $n\ge N$ 时 $I_n>I-\varepsilon$。
固定这个 $D_N$，取
\[
M=\sup\{\sqrt{x^2+y^2}:(x,y)\in D_N\}.
\]
当 $d(\Gamma)>M$ 时，$D_\Gamma\supset D_N$，因此
\[
\displaystyle\iint_{D_\Gamma}f(x,y)\,\dd x\,\dd y
\ge\displaystyle\iint_{D_N}f(x,y)\,\dd x\,\dd y
=I_N>I-\varepsilon.
\]
结合上面的上界，得到对所有 $d(\Gamma)>M$ 的截断曲线，都有
\[
I-\varepsilon<\displaystyle\iint_{D_\Gamma}f(x,y)\,\dd x\,\dd y\le I.
\]
所以
\[
\lim_{d(\Gamma)\to+\infty}\displaystyle\iint_{D_\Gamma}f(x,y)\,\dd x\,\dd y=I.
\]
即反常二重积分收敛。
\end{proof}

\begin{theorem}[比较判别法]\label{p2-improper-comparison}
设 $0\le f(x,y)\le g(x,y)$，且两函数在每个有界截域上可积。
\begin{enumerate}
\item 若 $\displaystyle\iint_Dg$ 收敛，则 $\displaystyle\iint_Df$ 收敛。
\item 若 $\displaystyle\iint_Df$ 发散，则 $\displaystyle\iint_Dg$ 发散。
\end{enumerate}
\end{theorem}
\begin{proof}
对上述穷尽序列，有
\[
0\le\displaystyle\iint_{D_n}f\le\displaystyle\iint_{D_n}g.
\]
若右侧有上界，左侧也有上界，第一条由引理得到。
若 $\displaystyle\iint_Df$ 发散而 $\displaystyle\iint_Dg$ 收敛，第一条又会推出 $\displaystyle\iint_Df$ 收敛，
产生矛盾，故第二条成立。
\end{proof}

\begin{theorem}[收敛与绝对收敛]\label{p2-improper-absolute}
设 $f$ 在 $D$ 的每个有界截域上可积，则
\[
\displaystyle\iint_Df(x,y)\,\dd x\,\dd y\text{ 收敛}
\quad\Longleftrightarrow\quad
\displaystyle\iint_D|f(x,y)|\,\dd x\,\dd y\text{ 收敛}.
\]
\end{theorem}
\begin{proof}
设
\[
f^+(x,y)=\begin{cases}f(x,y),&f(x,y)\ge0,\\0,&f(x,y)<0,\end{cases}
\qquad f^+(x,y)=\max(f(x,y),0),
\]
\[
f^-(x,y)=\begin{cases}-f(x,y),&f(x,y)\le0,\\0,&f(x,y)>0,\end{cases}
\qquad f^-(x,y)=\max(-f(x,y),0).
\]
两函数均非负，且
\[
|f|=f^++f^-,\qquad f=f^+-f^-.
\]

\textbf{充分性。}设 $|f|$ 在 $D$ 上可积。
由 $0\le f^\pm\le f^++f^-=|f|$ 及比较判别法，$f^+$、$f^-$ 均在 $D$ 上可积。
再由 $f=f^+-f^-$，得到 $f$ 在 $D$ 上可积。

\textbf{必要性。}设 $f$ 在 $D$ 上可积，而 $|f|$ 在 $D$ 上不可积。
由 $|f|=f^++f^-$，$f^+$、$f^-$ 的反常积分至少有一个发散。
不妨设 $\displaystyle\iint_Df^+$ 发散。那么对任意 $k>0$，可找到一条截断曲线 $\Gamma'$，使
\[
\displaystyle\iint_{D_{\Gamma'}}f^+(x,y)\,\dd x\,\dd y>k.
\]

任取截断曲线 $\Gamma_1$，记它截取的区域为 $D_1$。
再取截断曲线，使 $D_2\supset D_1$，且 $D_2$ 上的正部积分足够大。
如此逐步选取，得到一列曲线 $\{\Gamma_n\}$，满足
\[
D_1\subset D_2\subset\cdots\subset D_n\subset\cdots,\qquad
d(\Gamma_n)\longrightarrow+\infty,
\]
并且
\[
\displaystyle\iint_{D_{n+1}}f^+(x,y)\,\dd x\,\dd y
>2\displaystyle\iint_{D_n}|f(x,y)|\,\dd x\,\dd y+n.
\]
选取时可同时要求 $D_{n+1}$ 包含 $D_n$ 及一个足够大的圆内的区域；
由非负性，扩大区域不会减小正部积分，所以这些要求能够同时满足。
由 $f^+\le|f|$，有
\begin{align*}
\displaystyle\iint_{D_{n+1}\setminus D_n}f^+
&>2\displaystyle\iint_{D_n}|f|+n-\displaystyle\iint_{D_n}f^+\\
&\ge\displaystyle\iint_{D_n}|f|+n.
\end{align*}

对 $D_{n+1}\setminus D_n$ 作划分，分成有限个可求面积的小区域，
各小区域的面积记为 $\Delta\sigma_k$，$f^+$ 在其中的下确界记为 $m_k$。
它们组成 $f^+$ 的一个 Darboux 下和。
由可积性，可使下和与积分之差小于 $1$，故
\[
\sum_{k=1}^{\ell}m_k\Delta\sigma_k
>\displaystyle\iint_{D_{n+1}\setminus D_n}f^+-1
>\displaystyle\iint_{D_n}|f|+n-1.
\]

把这些小区域分成两类：第 (I) 类为 $m_k>0$ 的小区域，
第 (II) 类为 $m_k=0$ 的小区域。
在第 (I) 类上 $f^+>0$，所以 $f=f^+$；
第 (II) 类对下和的贡献为零。
记所有第 (I) 类小区域的并为 $P_n$，则
\begin{align*}
\sum_{k=1}^{\ell}m_k\Delta\sigma_k
&=\sum_{\text{第 (I) 类}}m_k\Delta\sigma_k\\
&\le\displaystyle\iint_{P_n}f^+
=\displaystyle\iint_{P_n}f.
\end{align*}

令 $E_n=D_n\cup P_n$，则
\begin{align*}
\displaystyle\iint_{E_n}f
&=\displaystyle\iint_{E_n}f^+-\displaystyle\iint_{E_n}f^-\\
&\ge\displaystyle\iint_{D_n}|f|+n-1-\displaystyle\iint_{E_n}f^-\\
&=\displaystyle\iint_{D_n}|f|+n-1-\displaystyle\iint_{D_n}f^-\\
&\ge n-1.
\end{align*}
这里因为 $P_n$ 上 $f^-=0$，所以第三行把 $E_n$ 上的负部积分换成了 $D_n$ 上的负部积分。

问题是 $E_n=D_n\cup P_n$ 不一定连通。
这时我们可以用一些狭小的“走廊”将其连通后得到区域 $\widetilde E_n$，
并且让这些“走廊”的总面积能充分小，使得
\[
\displaystyle\iint_{\widetilde E_n}f(x,y)\,\dd x\,\dd y>n-2,
\qquad n=1,2,\ldots.
\]
这里连接有限个小区域只需有限条区域内部的曲线；
在包含这些曲线的有界区域内 $f$ 有界，因而可把走廊取得使新增部分的 $|f|$ 积分小于 $1$。
$\widetilde E_n$ 包含 $D_n$，所以其截断也趋于无穷。
上述积分却无界，这说明 $f$ 在 $D$ 上不可积，与假设矛盾。
\end{proof}


\begin{example}[无穷远处的幂函数]
设 $a>0$，
\[
D=\{(x,y):a^2\le x^2+y^2<+\infty\},\qquad
r=\sqrt{x^2+y^2},\qquad f(x,y)=r^{-p}\quad(p>0).
\]
则 $\displaystyle\iint_Df$ 在 $p>2$ 时收敛，在 $p\le2$ 时发散。
\end{example}
\begin{solution}
取 $\Gamma_\rho:\ x^2+y^2=\rho^2$（$\rho>a$），相应截域为
\[
D_\rho=\{(x,y):a^2\le x^2+y^2\le\rho^2\}.
\]
用极坐标变换，
\begin{align*}
\displaystyle\iint_{D_\rho}f(x,y)\,\dd x\,\dd y
&=\int_0^{2\pi}\dd\theta\int_a^\rho r^{-p}r\,\dd r\\
&=2\pi\int_a^\rho r^{1-p}\,\dd r.
\end{align*}
当 $p\ne2$ 时，右边为 $2\pi(\rho^{2-p}-a^{2-p})/(2-p)$；
当 $p=2$ 时，为 $2\pi\log(\rho/a)$。
所以 $p>2$ 时极限为 $2\pi a^{2-p}/(p-2)$，$p\le2$ 时趋于 $+\infty$。
由非负函数的穷尽判据，得到所述结论。
\end{solution}

\begin{corollary}[Cauchy 判别法]
设 $D$ 用极坐标表示为
\[
D=\{(r,\theta):a\le r<+\infty,\ \alpha\le\theta\le\beta\},
\quad a>0,\quad 0\le\alpha<\beta\le2\pi.
\]
设 $f$ 在每个有界截域上可积。
\begin{enumerate}
\item 若存在 $M>0$，使 $|f(x,y)|\le M/r^p$，其中 $p>2$，则 $\displaystyle\iint_Df$ 收敛。
\item 若存在 $m>0$，使 $|f(x,y)|\ge m/r^p$，其中 $p\le2$，则 $\displaystyle\iint_Df$ 发散。
\end{enumerate}
\end{corollary}
\begin{proof}
在这个扇形区域上，幂函数截域积分为
\[
\displaystyle\iint_{D_\rho}r^{-p}\,\dd x\,\dd y
=(\beta-\alpha)\int_a^\rho r^{1-p}\,\dd r.
\]
由于 $\beta-\alpha>0$，其敛散性与上例相同。
第一条先由比较判别法得 $\displaystyle\iint_D|f|$ 收敛，再用绝对收敛定理。
第二条由比较判别法得 $\displaystyle\iint_D|f|$ 发散，
若 $\displaystyle\iint_Df$ 收敛，则绝对收敛定理又要求它绝对收敛，矛盾。
这里要求角度区间有正长度；零面积的退化情形不属于第二条的结论。
\end{proof}
一般地，我们有对于反常二重积分的累次积分公式与变量代换公式.
\begin{theorem}[化为累次反常积分]\label{p2-improper-iterated}
设 $f$ 在 $D=[a,+\infty)\times[c,+\infty)$ 上连续，且累次反常积分
\[
\int_a^{+\infty}\dd x\int_c^{+\infty}f(x,y)\,\dd y,
\qquad
\int_a^{+\infty}\dd x\int_c^{+\infty}|f(x,y)|\,\dd y
\]
均存在（内层积分对每个 $x$ 存在，所得函数的外层反常积分也存在）。
则 $f$ 在 $D$ 上可积，并且
\[
\displaystyle\iint_Df(x,y)\,\dd x\,\dd y
=\int_a^{+\infty}\dd x\int_c^{+\infty}f(x,y)\,\dd y.
\]
\end{theorem}


\begin{theorem}[反常二重积分的变量代换]
设 $T:D\to T(D)$ 是一一对应的连续可导坐标变换，
\[
x=x(u,v),\qquad y=y(u,v),\qquad
\frac{\partial(x,y)}{\partial(u,v)}\ne0.
\]
在相应截域穷尽各自区域、并避开可能的奇点的意义下，有
\[
\displaystyle\iint_{T(D)}f(x,y)\,\dd x\,\dd y
=\displaystyle\iint_Df\bigl(x(u,v),y(u,v)\bigr)
\left|\frac{\partial(x,y)}{\partial(u,v)}\right|\,\dd u\,\dd v.
\]
只要其中一边收敛，另一边也收敛，且两者相等。
\end{theorem}

\begin{example}[Poisson 积分]
计算 $\displaystyle\iint_{\R^2}e^{-(x^2+y^2)}\,\dd x\,\dd y$，并求
$\displaystyle\int_0^{+\infty}e^{-x^2}\,\dd x$。
\end{example}
\begin{solution}
利用极坐标变换 $x=r\cos\theta$、$y=r\sin\theta$，$\R^2$ 就变换为
\[
D=\{(r,\theta):0\le r<+\infty,\ 0\le\theta\le2\pi\}.
\]
因此利用变量代换法得
\begin{align*}
\displaystyle\iint_{\R^2}e^{-(x^2+y^2)}\,\dd x\,\dd y
&=\displaystyle\iint_D e^{-r^2}r\,\dd r\,\dd\theta\\
&=\int_0^{2\pi}\dd\theta\int_0^{+\infty}e^{-r^2}r\,\dd r\\
&=2\pi\left[-\frac12e^{-r^2}\right]_0^{+\infty}=\pi.
\end{align*}


又由于 $\R^2=(-\infty,+\infty)\times(-\infty,+\infty)$，所以利用化累次积分法得
\begin{align*}
\pi
&=\displaystyle\iint_{\R^2}e^{-(x^2+y^2)}\,\dd x\,\dd y\\
&=\int_{-\infty}^{+\infty}\dd x\int_{-\infty}^{+\infty}
 e^{-(x^2+y^2)}\,\dd y\\
&=\left(\int_{-\infty}^{+\infty}e^{-x^2}\,\dd x\right)
  \left(\int_{-\infty}^{+\infty}e^{-y^2}\,\dd y\right)=J^2.
\end{align*}
由 $J>0$ 得 $J=\sqrt\pi$。又因被积函数为偶函数，
\[
\int_0^{+\infty}e^{-x^2}\,\dd x=\frac{\sqrt\pi}{2}.
\]
最后这个积分称为 Poisson 积分，在概率统计等领域中有重要应用。
\end{solution}

\subsection*{无界函数的反常重积分}

设 $D$ 为有界区域，$P_0\in D$。
若 $f$ 在 $D\setminus\{P_0\}$ 上有定义，
但在 $P_0$ 的任意去心邻域内都无界，则称 $P_0$ 为 $f$ 的一个\textbf{奇点}。
设 $\gamma$ 为内部包含 $P_0$ 的零面积闭曲线，$\sigma$ 为它围成的区域。
总假设删去这个邻域后的通常积分
\[
\displaystyle\iint_{D\setminus\sigma}f(x,y)\,\dd x\,\dd y
\]
存在。这里删除的是奇点附近的一整块区域，不是只删除奇点本身。

\begin{definition}[单个奇点的反常二重积分]
记
\[
\rho(\gamma)=\sup\{|P-P_0|:P\in\gamma\}.
\]
若当 $\rho(\gamma)\to0$ 时，$\displaystyle\iint_{D\setminus\sigma}f$ 对所有这样的删域方式
都趋于同一个有限极限，则称反常二重积分收敛，并记
\[
\displaystyle\iint_Df(x,y)\,\dd x\,\dd y
=\lim_{\rho(\gamma)\to0}\displaystyle\iint_{D\setminus\sigma}f(x,y)\,\dd x\,\dd y.
\]
若该有限极限不存在，则称反常二重积分发散。
\end{definition}

如果函数 $f(x,y)$ 在区域 $D$ 上有奇线 $\Gamma_0$，
即 $f(x,y)$ 在 $D\setminus\Gamma_0$ 上有定义，但在任何包含曲线 $\Gamma_0$ 的区域上无界，
则可与上面的定义一样规定 $f(x,y)$ 在 $D$ 上的反常二重积分。

\textbf{奇线情形的定义}
若奇点不是单个点，而是 $D$ 内的一条零面积曲线 $\Gamma_0$，也可以按同样的方式定义。
取包含 $\Gamma_0$ 的闭邻域 $\sigma$，要求 $\Gamma_0$ 包含于 $\sigma$ 的内部，
且删域的边界分段光滑。用
\[
\rho(\sigma)=\sup_{P\in\sigma}\operatorname{dist}(P,\Gamma_0),\qquad
\operatorname{dist}(P,\Gamma_0)=\inf_{Q\in\Gamma_0}|P-Q|
\]
衡量邻域的宽度；这里点到集合的距离由右边的下确界定义。
当 $\rho(\sigma)\to0$ 时，若 $\displaystyle\iint_{D\setminus\sigma}f$ 对任意这样的邻域
都有同一个有限极限，就把它定义为 $f$ 在 $D$ 上的反常二重积分。
对有限个奇点或有限条奇异曲线，分别删去它们的邻域并令各邻域同时缩小，定义相同。



